% ===========================================================================
\part{Data}
% ===========================================================================

\chapter{The data model, as it is and as it should be}
\label{ch:datamodel}

\section{What the four migrations built}

Figure~\ref{fig:er-now} is the schema currently in the database, read back from Postgres rather
than from the migration files, so it reflects what actually exists.

\begin{diagram}{The current schema. The three relationships are all modelled as one-to-many with
the foreign key on the child; two of the three should not be. Every foreign-key column carries its
own sequence (see \S\ref{sec:serialfk}).}{fig:er-now}
\begin{tikzpicture}[node distance=8mm]
  \node[ent,minimum width=42mm] (pi) {\textbf{personal\_informations}\nodepart{two}
    \scriptsize id \textsc{pk}\\\scriptsize name, surname\\\scriptsize image\_url\\\scriptsize birth\_date};
  \node[ent,minimum width=42mm,right=22mm of pi] (ci) {\textbf{contact\_informations}\nodepart{two}
    \scriptsize id \textsc{pk}\\\scriptsize personal\_information\_id \textsc{fk}\\\scriptsize github, email\\\scriptsize instagram, linked\_in};
  \draw[flow] (pi) -- node[lbl,above]{1} node[lbl,below]{N} (ci);

  \node[ent,minimum width=42mm,below=13mm of pi] (je) {\textbf{job\_experiences}\nodepart{two}
    \scriptsize id \textsc{pk}\\\scriptsize date\_from, date\_to\\\scriptsize job\_title\\\scriptsize accomplishments\\\scriptsize responsabilities};
  \node[ent,minimum width=42mm,right=22mm of je] (ji) {\textbf{job\_institutions}\nodepart{two}
    \scriptsize id \textsc{pk}\\\scriptsize job\_experience\_id \textsc{fk}\\\scriptsize name\\\scriptsize url};
  \draw[flow] (je) -- node[lbl,above]{1} node[lbl,below]{N} (ji);

  \node[ent,minimum width=42mm,below=13mm of je] (po) {\textbf{portfolio\_items}\nodepart{two}
    \scriptsize id \textsc{pk}\\\scriptsize title \textsc{unique}\\\scriptsize description\\\scriptsize public\\\scriptsize public\_url};
  \node[ent,minimum width=42mm,right=22mm of po] (tg) {\textbf{tags}\nodepart{two}
    \scriptsize id \textsc{pk}\\\scriptsize portfolio\_item\_id \textsc{fk}\\\scriptsize value};
  \draw[flow] (po) -- node[lbl,above]{1} node[lbl,below]{N} (tg);

  \node[ent,minimum width=42mm,below=13mm of po] (pu) {\textbf{publication\_items}\nodepart{two}
    \scriptsize id \textsc{pk}\\\scriptsize title, journal\\\scriptsize abs\\\scriptsize year\\\scriptsize link};
  \node[tag,right=22mm of pu,align=left,text width=40mm] {no relationships\\--- correctly so, for now};
\end{tikzpicture}
\end{diagram}

The structure is sound in outline: four content types, each with the fields a personal site
actually needs. What follows is not a rejection of it. It is the set of decisions that are cheap to
change now, while every table is empty, and expensive to change once there is content and a
deployed site.

\section{Three modelling decisions to revisit}

\subsection{\code{contact\_informations} is one-to-one pretending to be one-to-many}

There is one of you. There will be one row in \code{personal\_informations} and one row in
\code{contact\_informations}. The schema permits four contact rows for the same person, and
nothing stops it, so every query that reads contact details has to decide what to do about the
possibility --- \code{get\_result} will error on multiple rows, \code{get\_results} forces the
caller to handle a \code{Vec} that should always have length one.

There is also a second, quieter problem: the columns are the platforms. \code{github},
\code{email}, \code{instagram}, \code{linked\_in}. Adding Mastodon means a migration, a schema
regeneration, a model change, a DTO change and a component change. Removing Instagram means the
same in reverse, or a dead column forever.

\begin{pattern}[Make contact links rows, not columns]
\begin{lstlisting}[style=sql]
CREATE TABLE contact_links (
  id          integer GENERATED ALWAYS AS IDENTITY PRIMARY KEY,
  person_id   integer NOT NULL REFERENCES people(id) ON DELETE CASCADE,
  kind        text    NOT NULL,           -- 'github' | 'email' | 'linkedin' | ...
  label       text    NOT NULL,           -- what to show: "@egonik"
  url         text    NOT NULL,
  position    integer NOT NULL DEFAULT 0, -- display order
  UNIQUE (person_id, kind)
);
CREATE INDEX contact_links_person_idx ON contact_links (person_id, position);
\end{lstlisting}
Now a new platform is one \code{INSERT}. The component iterates rows and renders an icon chosen
from \code{kind}, and adding a platform never touches Rust.
\end{pattern}

For the person table itself, if there is genuinely only ever one row, say so in the schema rather
than hoping:

\begin{lstlisting}[style=sql]
-- a singleton table: exactly one row, enforced
CREATE TABLE people (
  id         integer PRIMARY KEY GENERATED ALWAYS AS IDENTITY,
  ...
  CONSTRAINT people_singleton CHECK (id = 1)
);
\end{lstlisting}

\subsection{\code{job\_institutions} points the wrong way}

This is the one worth changing before there is data, because the fix touches a foreign key.

As modelled, one \code{job\_experience} has many \code{job\_institutions}, and each institution
belongs to exactly one experience. Read that as a sentence about the world: \emph{one job was held
at several companies, and each company hosted exactly one job.} That is backwards. A job is held at
one institution, and an institution can host several jobs --- which is what happens the second time
you are promoted, or return to a former employer.

The practical consequence is duplication. Work at Snappler twice and you get two
\code{job\_institutions} rows both named ``Snappler'', both with the same URL, and no way to know
they are the same organisation. Fix the URL and you fix it once, in one of two places.

\begin{pattern}[Institutions are their own entity]
\begin{lstlisting}[style=sql]
CREATE TABLE institutions (
  id   integer GENERATED ALWAYS AS IDENTITY PRIMARY KEY,
  name text NOT NULL UNIQUE,
  url  text
);

ALTER TABLE job_experiences
  ADD COLUMN institution_id integer NOT NULL REFERENCES institutions(id) ON DELETE RESTRICT;
CREATE INDEX job_experiences_institution_idx ON job_experiences (institution_id);
\end{lstlisting}
\code{ON DELETE RESTRICT} rather than \code{CASCADE}: deleting an institution should not silently
delete your employment history. The right answer is to refuse.
\end{pattern}

\subsection{\code{tags} should be many-to-many}

\code{tags(id, portfolio\_item\_id, value)} stores the tag \emph{text} once per item. Tag three
projects with ``rust'' and the string ``rust'' exists three times. Consequences: no canonical list
of tags to render a filter UI from, no way to rename a tag in one place, and typos
(``leptos''/``Leptos'') that silently become two different tags.

\begin{pattern}[The standard join-table shape]
\begin{lstlisting}[style=sql]
CREATE TABLE tags (
  id    integer GENERATED ALWAYS AS IDENTITY PRIMARY KEY,
  slug  text NOT NULL UNIQUE,             -- 'rust'      -- for URLs
  label text NOT NULL                     -- 'Rust'      -- for display
);

CREATE TABLE portfolio_item_tags (
  portfolio_item_id integer NOT NULL REFERENCES portfolio_items(id) ON DELETE CASCADE,
  tag_id            integer NOT NULL REFERENCES tags(id)            ON DELETE CASCADE,
  PRIMARY KEY (portfolio_item_id, tag_id)
);
CREATE INDEX portfolio_item_tags_tag_idx ON portfolio_item_tags (tag_id);
\end{lstlisting}
The composite primary key makes duplicate tagging impossible, and the extra index makes
``all projects tagged rust'' fast in the other direction. Diesel handles this shape well ---
\code{belonging\_to} plus a grouped load, which is the standard pattern in its guides.
\end{pattern}

\section{Column-level corrections}

\begin{table}[htbp]
\centering
\caption{Column changes worth making while the tables are empty.}
\label{tab:columns}
\small
\begin{tabular}{@{}L{0.30\textwidth}L{0.30\textwidth}L{0.34\textwidth}@{}}
\toprule
\textbf{Now} & \textbf{Should be} & \textbf{Why} \\
\midrule
\code{responsabilities} & \code{responsibilities} & Misspelling. It is about to become a column
name, a struct field, a JSON key and part of a public API. Free to fix today. \dnum{15} \\
\code{abs} \code{TEXT} & \code{abstract\_text} & \code{abs} reads as the SQL \code{ABS()} function.
Not reserved, so it works --- but nobody will guess what it holds. \dnum{15} \\
\code{accomplishments VARCHAR(255)} & child table, or \code{TEXT} & 255 characters is one sentence.
Accomplishments are a list. A \code{job\_highlights} child table with \code{(kind, position, body)}
lets you render bullets. \dnum{16} \\
\code{responsabilities VARCHAR(255)} & same & As above. \\
\code{public BOOLEAN} & \code{visibility} enum, or keep & Two-state today, three tomorrow
(draft / unlisted / public). An enum type or a \code{text} + \code{CHECK} avoids a second boolean. \\
\code{year INTEGER} & keep & Correct. A publication year is a year, not a date. \\
--- & \code{slug TEXT NOT NULL UNIQUE} & On every content table. You want
\code{/portfolio/egonik-site}, not \code{/portfolio/7}. Adding this later means either breaking
URLs or maintaining both. \\
--- & \code{position INTEGER NOT NULL DEFAULT 0} & Every list on the site has an order you care
about, and it is never ``by primary key''. \\
--- & \code{created\_at}, \code{updated\_at} \code{TIMESTAMPTZ} & Cheap now, impossible to
backfill later. Diesel's initial-setup migration already installed
\code{diesel\_manage\_updated\_at()} for exactly this and it is unused. \\
\bottomrule
\end{tabular}
\end{table}

\begin{pattern}[Use \code{GENERATED ALWAYS AS IDENTITY}, not \code{SERIAL}]
Every table currently uses \code{SERIAL} for its primary key, which is the older mechanism:
\code{SERIAL} is shorthand for ``integer, plus a sequence, plus a default'', and the sequence is
only loosely attached to the column. Identity columns are the SQL-standard replacement, are owned
properly by the column, and --- crucially for the next section --- cannot be applied by accident to
a column that should not have a default.
\end{pattern}

\section{The \code{SERIAL} foreign key, which is a live bug}
\label{sec:serialfk}

All three foreign-key columns are declared \code{SERIAL}:

\begin{lstlisting}[style=sql,caption={From the three \file{up.sql} files.}]
portfolio_item_id       serial not null,   -- tags
job_experience_id       serial,            -- job_institutions
personal_information_id serial not null,   -- contact_informations
\end{lstlisting}

The intent was ``an integer that references another table''. What \code{SERIAL} actually does is
create a \emph{dedicated sequence for this column} and set it as the column default. The
consequence is that a foreign key that should always be supplied explicitly instead has a default
value that counts up on its own, entirely independently of the parent table.

\begin{verified}[Verified: the real column definitions in the live database]
\begin{lstlisting}[style=out]
$ psql -d my-site -c '\d tags'
      Column       |  Type   | Nullable |                     Default
-------------------+---------+----------+-------------------------------------------------
 id                | integer | not null | nextval('tags_id_seq'::regclass)
 portfolio_item_id | integer | not null | nextval('tags_portfolio_item_id_seq'::regclass)
 value             | varchar |          |
\end{lstlisting}
There is a \code{tags\_portfolio\_item\_id\_seq}. A sequence, on a foreign key.
\end{verified}

That is not merely untidy. It converts a mistake that Postgres would have caught into silent
data corruption:

\begin{verified}[Verified: silent mis-attachment, then a confusing error]
\begin{lstlisting}[style=out]
INSERT INTO portfolio_items (title, description, public) VALUES ('Project A','desc',true);
INSERT INTO portfolio_items (title, description, public) VALUES ('Project B','desc',true);

-- the author forgets the foreign key. There is a DEFAULT, so this SUCCEEDS:
INSERT INTO tags (value) VALUES ('rust');      -- INSERT 0 1
INSERT INTO tags (value) VALUES ('leptos');    -- INSERT 0 1

SELECT t.value, p.title AS silently_attached_to FROM tags t
  JOIN portfolio_items p ON p.id = t.portfolio_item_id;

  value  | silently_attached_to
 --------+----------------------
  rust   | Project A                <-- wrong, and no error was raised
  leptos | Project B                <-- wrong

INSERT INTO tags (value) VALUES ('this-one-explodes');
ERROR:  insert or update on table "tags" violates foreign key constraint
DETAIL:  Key (portfolio_item_id)=(3) is not present in table "portfolio_items".
\end{lstlisting}
\end{verified}

Read that sequence of events carefully, because it is the worst shape a bug can have. The first two
inserts are wrong and succeed. The third fails --- with an error that points at the foreign key,
which is where you will look, but the actual cause is a default on a column three lines up in a
migration written two days earlier. Had the column been a plain \code{integer NOT NULL}, the very
first insert would have failed immediately with \code{null value in column "portfolio\_item\_id"
violates not-null constraint}, which is a message that leads directly to the mistake.

This is defect \dnum{1}, and it is the highest-priority item in this document. The fix is
mechanical:

\begin{lstlisting}[style=sql,caption={A migration that removes the three stray defaults and drops the sequences.}]
ALTER TABLE tags                 ALTER COLUMN portfolio_item_id       DROP DEFAULT;
ALTER TABLE job_institutions     ALTER COLUMN job_experience_id       DROP DEFAULT;
ALTER TABLE contact_informations ALTER COLUMN personal_information_id DROP DEFAULT;

DROP SEQUENCE IF EXISTS tags_portfolio_item_id_seq;
DROP SEQUENCE IF EXISTS job_institutions_job_experience_id_seq;
DROP SEQUENCE IF EXISTS contact_informations_personal_information_id_seq;

-- job_institutions.job_experience_id is also nullable, which it should not be
ALTER TABLE job_institutions ALTER COLUMN job_experience_id SET NOT NULL;
\end{lstlisting}

Because all four tables are empty, the honest alternative is better: rewrite the four
\file{up.sql} files correctly, drop the database, and re-run from scratch.
Section~\ref{sec:rewrite} explains when that is legitimate.

\section{Missing constraints and indexes}

Postgres will enforce whatever you tell it to and nothing more. The current schema tells it very
little. This checklist is worth running against every future migration.

\begin{table}[htbp]
\centering
\caption{Constraint checklist, with the current state of each item.}
\label{tab:constraints}
\small
\begin{tabular}{@{}L{0.30\textwidth}cL{0.52\textwidth}@{}}
\toprule
\textbf{Constraint} & \textbf{Now} & \textbf{Why it earns its place} \\
\midrule
Index on every FK column & \textcolor{crimson}{none} & Postgres indexes the \emph{referenced} key
automatically but never the referencing column. Be honest about the reason: at a few dozen rows
Postgres will seq-scan whichever way you index it, and should. The argument is that
\code{ON DELETE CASCADE} takes row locks on the child scan, and that the index costs nothing. \\
\code{ON DELETE} on every FK & \textcolor{crimson}{none} & The default is \code{NO ACTION}, so
deleting a portfolio item with tags fails with a foreign-key error. Decide deliberately:
\code{CASCADE} for tags, \code{RESTRICT} for institutions. \\
\code{CHECK (date\_to IS NULL OR date\_to >= date\_from)} & \textcolor{crimson}{no} & A job that
ended before it began will render as a negative duration. One line prevents it forever. \\
\code{CHECK (year BETWEEN 1900 AND 2100)} & \textcolor{crimson}{no} & Catches a typo'd year at
insert time rather than in a rendered page. \\
\code{UNIQUE} on \code{slug} & n/a & Once slugs exist, this is what makes a URL a key. \\
\code{NOT NULL} where meant & partial & \code{job\_institutions.job\_experience\_id} is nullable ---
an institution attached to no job. Almost certainly unintended. \\
\code{title UNIQUE} on \code{portfolio\_items} & \textcolor{moss}{yes} & Already there. Good ---
though once slugs exist, the slug is the better uniqueness key and the title constraint can relax. \\
\bottomrule
\end{tabular}
\end{table}

\section{The target schema, decided}
\label{sec:targetschema}

The sections above argue for individual changes in isolation, which is how you evaluate them and a
bad way to apply them. Here is the whole thing decided at once, so that M0 in
Chapter~\ref{ch:roadmap} is a single coherent migration rather than a dozen negotiations.

\begin{diagram}{The target schema. Every relationship has an explicit cardinality and referential
action; every content table has a \code{slug}, a \code{position} and timestamps. Compare with
Figure~\ref{fig:er-now}.}{fig:er-target}
\begin{tikzpicture}[node distance=7mm]
  % --- person cluster ---
  \node[ent,minimum width=40mm] (pr) {\textbf{profile}\nodepart{two}
    \scriptsize id \textsc{pk} \textsc{check} $=1$\\\scriptsize name, surname\\\scriptsize image\_url, birth\_date\\\scriptsize created\_at, updated\_at};
  \node[ent,minimum width=40mm,right=20mm of pr] (cl) {\textbf{contact\_links}\nodepart{two}
    \scriptsize id \textsc{pk}\\\scriptsize profile\_id \textsc{fk} \textsc{cascade}\\\scriptsize kind, label, url\\\scriptsize position\\\scriptsize \textsc{unique}(profile\_id, kind)};
  \draw[flow] (pr) -- node[lbl,above]{1} node[lbl,below]{N} (cl);

  % --- experience cluster ---
  \node[ent,minimum width=40mm,below=11mm of pr] (inst) {\textbf{institutions}\nodepart{two}
    \scriptsize id \textsc{pk}\\\scriptsize name \textsc{unique}\\\scriptsize url};
  \node[ent,minimum width=40mm,right=20mm of inst] (je) {\textbf{job\_experiences}\nodepart{two}
    \scriptsize id \textsc{pk}\\\scriptsize institution\_id \textsc{fk} \textsc{restrict}\\\scriptsize job\_title\\\scriptsize date\_from, date\_to?\\\scriptsize \textsc{check} to $\geq$ from};
  \node[ent,minimum width=40mm,right=20mm of je] (jh) {\textbf{job\_highlights}\nodepart{two}
    \scriptsize id \textsc{pk}\\\scriptsize job\_experience\_id \textsc{fk} \textsc{cascade}\\\scriptsize kind, position\\\scriptsize body \textsc{text}};
  \draw[flow] (inst) -- node[lbl,above]{1} node[lbl,below]{N} (je);
  \draw[flow] (je) -- node[lbl,above]{1} node[lbl,below]{N} (jh);

  % --- portfolio cluster ---
  \node[ent,minimum width=40mm,below=11mm of inst] (po) {\textbf{portfolio\_items}\nodepart{two}
    \scriptsize id \textsc{pk}\\\scriptsize slug \textsc{unique}\\\scriptsize title, summary\\\scriptsize description \textsc{text}\\\scriptsize is\_published\\\scriptsize live\_url?, repo\_url?};
  \node[ent,minimum width=40mm,right=20mm of po] (pit) {\textbf{portfolio\_item\_tags}\nodepart{two}
    \scriptsize portfolio\_item\_id \textsc{fk}\\\scriptsize tag\_id \textsc{fk}\\\scriptsize \textsc{pk}(both) \textsc{cascade}};
  \node[ent,minimum width=40mm,right=20mm of pit] (tg) {\textbf{tags}\nodepart{two}
    \scriptsize id \textsc{pk}\\\scriptsize slug \textsc{unique}\\\scriptsize label};
  \draw[flow] (po) -- node[lbl,above]{1} node[lbl,below]{N} (pit);
  \draw[flow] (tg) -- node[lbl,above]{1} node[lbl,below]{N} (pit);
  \node[lbl,below=1mm of pit] {effectively M:N};

  % --- publications ---
  \node[ent,minimum width=40mm,below=11mm of po] (pu) {\textbf{publications}\nodepart{two}
    \scriptsize id \textsc{pk}\\\scriptsize slug \textsc{unique}\\\scriptsize title, journal\\\scriptsize abstract\_text \textsc{text}\\\scriptsize year \textsc{check}, doi?, url};
  \node[tag,right=20mm of pu,align=left,text width=52mm] {standalone. Add an \code{authors} child
    table only when a publication has co-authors you want to link.};

  \node[tag,below=8mm of pu,align=left,text width=100mm] {\emph{Every table also carries}
    \code{created\_at TIMESTAMPTZ NOT NULL DEFAULT now()} \emph{and} \code{updated\_at}, \emph{and
    every primary key is} \code{GENERATED ALWAYS AS IDENTITY}. \emph{Every FK column has an index.}};
\end{tikzpicture}
\end{diagram}

Five renames are folded in, and each is here because the current name is either wrong or ambiguous:

\begin{itemize}
  \item \code{personal\_informations} $\rightarrow$ \code{profile}. The old name is ungrammatical
    (``informations'') and plural for a table that holds one row.
  \item \code{contact\_informations} $\rightarrow$ \code{contact\_links}, restructured from columns
    to rows.
  \item \code{publication\_items} $\rightarrow$ \code{publications}. The \code{\_item} suffix carries
    no information; \code{portfolio\_items} keeps it because a portfolio \emph{is} a collection of
    items.
  \item \code{abs} $\rightarrow$ \code{abstract\_text}; \code{responsabilities} $\rightarrow$ a
    \code{kind} value in \code{job\_highlights}.
  \item \code{public} / \code{public\_url} $\rightarrow$ \code{is\_published} / \code{live\_url}.
    Those two columns share a prefix and mean unrelated things --- one is a visibility flag, the
    other is a hyperlink --- and the shared prefix invites exactly the wrong mental model.
\end{itemize}

\begin{pattern}[Get here in one migration, not twelve]
Because every table is empty (Chapter~\ref{ch:state}), this is not a migration sequence --- it is a
rewrite of the four existing \file{up.sql} files plus \code{diesel database reset}. Do not write
twelve \code{ALTER TABLE} migrations to reach a schema that has never held a row.
\S\ref{sec:rewrite} sets out when that shortcut is legitimate, and why the door closes the day you
deploy.
\end{pattern}

% ===========================================================================
\chapter{Migration discipline}
\label{ch:migrations}

\section{The rule}

\begin{keyidea}
A migration is not ``the SQL that got the schema into the state I wanted''. It is a \emph{pair} of
scripts, and the pair is only correct when \code{up} followed by \code{down} returns the database
to exactly where it started. If \code{down} does not work, you do not have a migration --- you have
a one-way change with a file next to it.
\end{keyidea}

Right now \code{down} does not work. Not for one migration --- for three of the four, and in a way
that makes the whole chain unrevertible.

\section{Verified: the chain cannot be rolled back}

\begin{verified}[Verified on a scratch database]
A fresh database was created, all migrations applied cleanly, and then reverted one at a time:
\begin{lstlisting}[style=out]
$ diesel migration run
Running migration 2026-08-02-211314-0000_create_portfolio_item
Running migration 2026-08-03-231439-0000_create_publications_item
Running migration 2026-08-03-232042-0000_create_job_experience
Running migration 2026-08-03-233332-0000_create_personal_information

$ diesel migration revert
Rolling back migration 2026-08-03-233332-0000_create_personal_information
ERROR diesel::database: Failed to execute query
  query="drop table personal_informations;\ndrop table contact_informations;\n"
  err=DatabaseError(Unknown,
      "cannot drop table personal_informations because other objects depend on it")
Failed to run migrations: Failed to run 2026-08-03-233332-0000_create_personal_information
  with: cannot drop table personal_informations because other objects depend on it
\end{lstlisting}
Running it four more times produces the identical error. The chain is stuck at the top: because the
most recent migration cannot be reverted, \emph{none} of the earlier ones can be reached.
\end{verified}

\subsection{Cause 1: drops in the wrong order}

\begin{lstlisting}[style=sql,caption={\file{2026-08-03-233332-0000\_create\_personal\_information/down.sql}}]
drop table personal_informations;      -- <-- the PARENT, dropped first
drop table contact_informations;       -- <-- the CHILD, which still references it
\end{lstlisting}

\code{contact\_informations} holds a foreign key into \code{personal\_informations}, so the parent
cannot be dropped while the child exists. \code{DROP} order must be the reverse of \code{CREATE}
order: children first, parents last. The \code{create\_job\_experience} migration has the identical
bug. Interestingly, \code{create\_portfolio\_item} gets it right --- \code{DROP table tags} before
\code{DROP TABLE portfolio\_items} --- which shows the correct order was understood once and then
lost. Defect \dnum{3}.

\subsection{Cause 2: a typo'd table name}

\begin{lstlisting}[style=sql,caption={\file{2026-08-03-231439-0000\_create\_publications\_item/down.sql}}]
drop table publication_itens
\end{lstlisting}

\begin{verified}
\begin{lstlisting}[style=out]
$ psql -d scratch -f migrations/..._create_publications_item/down.sql
ERROR:  table "publication_itens" does not exist
\end{lstlisting}
\end{verified}

\code{itens} is the Portuguese plural of \emph{item}; the table is \code{publication\_items}. A
one-character slip that makes the migration unrevertible, and that no amount of testing the
\code{up} direction would ever surface. Defect \dnum{4}.

\subsection{The corrected files}

\begin{lstlisting}[style=sql,caption={All three \file{down.sql} files, fixed. \code{IF EXISTS} makes a
partially-applied revert re-runnable.}]
-- create_publications_item/down.sql
DROP TABLE IF EXISTS publication_items;

-- create_job_experience/down.sql
DROP TABLE IF EXISTS job_institutions;      -- child first
DROP TABLE IF EXISTS job_experiences;       -- then parent

-- create_personal_information/down.sql
DROP TABLE IF EXISTS contact_informations;  -- child first
DROP TABLE IF EXISTS personal_informations; -- then parent
\end{lstlisting}

\section{The habit that would have caught all of this}

\begin{pattern}[Never write a migration without running \code{redo}]
\code{diesel migration redo} reverts the most recent migration and immediately re-applies it. It is
one command, it takes a second, and it exercises the \code{down} path that you will otherwise
discover is broken on the day you actually need it.
\begin{lstlisting}[style=sh]
diesel migration generate create_thing     # writes up.sql and down.sql
$EDITOR migrations/*_create_thing/{up,down}.sql
diesel migration run                       # apply
diesel migration redo                      # <-- the step that was skipped
diesel print-schema > src/schema.rs        # regenerate, then commit both together
\end{lstlisting}
Make the \code{redo} step non-negotiable. Every defect in this chapter would have been caught by it
within seconds of being written.
\end{pattern}

\begin{pitfall}[\code{redo} on a migration whose \code{down} is broken leaves you stuck]
Diesel runs each migration in a transaction, so a failed revert rolls back and the database is
unharmed. But the migration stays applied and \code{revert} will keep failing, so you cannot reach
any earlier migration until you fix the \code{down.sql}. That is the state this repository is in
right now.
\end{pitfall}

\section{When to rewrite history instead}
\label{sec:rewrite}

Normally, editing an already-applied migration is forbidden --- other people's databases and your
production database have recorded that it ran, and changing it silently desynchronises them.

That reasoning does not apply here, and it is worth saying so explicitly rather than following the
rule out of habit. There is exactly one database, on your machine, with zero rows in every table,
and one commit in the repository. Nothing has been deployed. Nobody else has cloned it.

\begin{pattern}[Reset now; never again after the first deploy]
\begin{lstlisting}[style=sh]
# 1. fix all four up.sql (SERIAL -> integer on FKs, constraints, indexes,
#    naming, created_at/updated_at) and all four down.sql (order, typo)
# 2. throw the database away and rebuild it from the fixed migrations
diesel database reset          # drop + create + run all migrations
# 3. prove the pair is symmetric, for every migration, not just the last
diesel migration redo -a       # -a = redo all
# 4. regenerate the schema and commit migrations + schema.rs together
diesel print-schema > src/schema.rs
\end{lstlisting}
After the site is deployed once, this door closes. From then on every schema change is a new
migration, forward-only, and the \code{down} script exists for your development loop rather than
for production rollback.
\end{pattern}

\begin{pitfall}[\code{schema.rs} is generated --- never hand-edit it]
\file{src/schema.rs} carries the header
\code{// @generated automatically by Diesel CLI}. It is the compiler's model of your tables, and
if it drifts from the database you get type errors that describe a schema nobody has. Regenerate it
after every migration, in the same commit as the migration, so that a single \code{git show}
tells the whole story of a schema change.
\end{pitfall}

\section{Housekeeping in the migrations directory}

Three small things, all defect-register entries:

\begin{description}
  \item[\file{migrations/.diesel\_lock} is untracked and unignored.] It is an ephemeral lock file
    created while the CLI runs. Add it to \file{.gitignore}. \dnum{21}

  \item[The timestamps carry a \code{-0000} suffix.] Diesel's own format is
    \code{YYYY-MM-DD-HHMMSS\_name}; yours are \code{2026-08-02-211314-0000\_create\_portfolio\_item}.
    Harmless --- ordering still works --- but let \code{diesel migration generate} name the
    directories so they stay consistent. \dnum{22}

  \item[Names should be plural and consistent.] \code{create\_portfolio\_item} creates
    \code{portfolio\_items} \emph{and} \code{tags}; \code{create\_publications\_item} mixes plural
    and singular in one name. Pick ``one migration named after what it creates'' and stick to it. \dnum{22}
\end{description}

% ===========================================================================
\chapter{Content: seeding, and who writes it}
\label{ch:seed}

\section{Seed before you render}

Every table is empty. That matters more than it sounds, because \textbf{an empty list and a broken
query render identically}. If you build the publications page against an empty table you cannot
tell whether it works, and the first time you insert a row you will be debugging the page and the
query at once.

\begin{pattern}[A seed migration, or a seed binary]
Two reasonable mechanisms, and the choice depends on whether the content \emph{is} the data or
merely stands in for it.

\textbf{For real content --- your actual CV --- use an ordinary migration.} It is versioned,
reviewed, applied identically everywhere, and rolls back with the schema. Your employment history
changes about once a year; a migration per change is entirely proportionate.

\textbf{For development fixtures --- ten fake portfolio items to test pagination --- use a small
binary} at \file{src/bin/seed.rs}, run with \code{cargo run --bin seed --features ssr}. Keep it out
of the migration chain so it never runs in production, and make it idempotent so you can run it
twice.
\end{pattern}

\begin{lstlisting}[style=sql,caption={A seed migration, with explicit foreign keys --- which
is what defect \dnum{1} currently lets you forget.}]
INSERT INTO people (id, name, surname, image_url, birth_date)
VALUES (1, 'Eduardo', 'Gonik', '/assets/me.jpg', '1990-01-01');

INSERT INTO contact_links (person_id, kind, label, url, position) VALUES
  (1, 'github',   '@egonik',        'https://github.com/egonik',            0),
  (1, 'email',    'ia@snappler.com','mailto:ia@snappler.com',               1),
  (1, 'linkedin', 'Eduardo Gonik',  'https://linkedin.com/in/...',          2);
\end{lstlisting}

\section{The authoring question you have not answered yet}

There is a decision waiting here that is easy to defer past the point where it is cheap. Every
repository in the codebase is read-only: \code{get\_all} and \code{get\_by\_id}. So how does content
get in?

Three answers, in increasing order of cost:

\begin{description}
  \item[SQL migrations, forever.] Content lives in the repository as SQL. Deploying is how you
    publish. No admin UI, no authentication, no attack surface. For a personal site updated a few
    times a year this is not a compromise --- it is arguably the correct answer, and it is what a
    static-site generator would give you with more machinery.

  \item[Markdown files in the repository, read at startup.] Content lives in
    \file{content/*.md} and is loaded into the database (or straight into memory) on boot. You get
    to write prose in an editor instead of inside a SQL string literal, and Postgres holds only
    what benefits from being queryable. A good fit if the site grows a blog.

  \item[An admin UI behind authentication.] Now you need write repositories, \code{Action}s and
    \code{<ActionForm>}s, session management, password hashing, CSRF considerations, and a
    permanent obligation to keep all of it patched. Worth it only if you will genuinely edit
    content from a phone.
\end{description}

\begin{pattern}[Choose the first, and revisit only when it hurts]
Start with SQL migrations. The reason is not laziness --- it is that server functions are public
HTTP endpoints. The moment a write path exists, it exists for everyone on the internet unless every
one of those functions checks authorisation on every call. That is a real, ongoing obligation, and
it is a strange one to take on for a site whose content changes annually.

The decision that matters is not which of the three you pick today. It is that you pick
\emph{deliberately}, and that if you defer, you defer knowing you have deferred --- rather than
discovering in six months that there is no way to add a project without hand-writing SQL.
\end{pattern}
